\documentclass[9pt]{article}
\usepackage[letterpaper,left=6.9cm,right=1.9cm,top=3.1cm,bottom=2.3cm]{geometry} % Márgenes y demás


\input{Styles_ACCEFN.sty}  %%Llama el archivo de los estilos


\setlength{\topskip}{0.5cm}
\setlength{\parindent}{0mm}
%\vspace{0.3cm} %Espacio reservado para la revista
{\title{\fontsize{14}{14}\textbf{Título del artículo en español, negrita, 14pt} \\[0.2cm]
\textcolor{gray}{\textbf{Título del artículo en Inglés, color gris, en negrita, 14pt}\\[0.2cm]}}
}

\author[1,*]{\fontsize{9}{9} \textbf{Autor 1}}
\author[2]{\textbf{ Autor 2}}
\author[1,3]{\textbf{Autor 3}}


\affil[1]{\fontsize{8}{8}Filiación Autor 1, 8pt}
\affil[2]{Filiación Autor 2}
\affil[3]{Filiación Autor 3.}
\date{}                     %% if you don't need date to appear
\setcounter{Maxaffil}{0}



%%%%%% Inicio del documento
\begin{document}

\maketitle
\section*{Resumen}
\justifying
Resumen en español e inglés que no supere 250 palabras. \lipsum[22] $S^{\alpha}_{\mu v}$ , Wisi enim ad minim veniam, quis nostrud
exerci tation ullamcorper suscipit lobortis nisl ut aliquip ex ea commodo consequat. Duis autem vel eum iriure dolor in hendrerit in vulputate velit esse molestie consequat.\\[0.2cm]
\textbf{Palabras clave:} Español; Inglés; Que no superen en número a seis; Separadas por punto y coma.


\section*{Abstract}

\justifying
Resumen en español e inglés que no supere 250 palabras. \lipsum[22]
$S^{\alpha}_{\mu v}$ , Wisi enim ad minim veniam, quis nostrud
exerci tation ullamcorper suscipit lobortis nisl ut aliquip ex ea commodo consequat. Duis autem vel eum iriure dolor in hendrerit in vulputate velit esse molestie consequat.\\[0.2cm]
\textbf{Keywords:} Español; Inglés; Que no superen en número a seis; Separadas por punto y coma.


\section*{Introduction}

Se debe mostrar el marco teórico en el cual se inscribe el tema de la investigación, propósito del artículo y justificación del estudio indicando, además, la delimitación de la pregunta
problema\cite{ostrogradski2007memoires}, \lipsum[20] \cite{woodard2007avoiding}; \lipsum[5] \cite{nicolis2009galileon}. 

\section*{Materiales y métodos}

Se describirán los sujetos de observación o experimentación, aparatos o dispositivos y metodologías empleadas en el estudio. La descripción de la metodología debe tener referencias de los métodos usados. Análisis estadístico: señalar los métodos estadísticos utilizados y la manera de analizar los datos.

\begin{equation}
    {\displaystyle \mathbf {\hat {H}} \left|\psi (t)\right\rangle =i\hbar {\partial  \over \partial t}\left|\psi (t)\right\rangle }
\end{equation}

siendo $F_{\mu v} \equiv \partial_{\mu}A_{v}-\partial_{\mu}A_{v}$ con $A_{\mu}$ siendo el campo vectorial masivo y $M$ alibium caret
forlmui equnetemeociemm. uis autem vel eum iriure dolor in hendrerit in vulputate velit esse molestie consequat, vel illum dolore eu feugiat nulla facilisis at vero eros et accumsan t wisi enim ad minim veniam, quis nostrud exerci tation ullamcorper suscipit lobortis nisl ut aliquip ex ea commodo consequat

\begin{equation}
    K = mx+b \dots
\end{equation}

\subsection*{Suntítulo de sección}

\lipsum[15]


{\setlength{\parindent}{6mm} Nostrud oui masonim exerci tation ullamcorper suscipit lobortis nisl ut ali commodo \cite{burton1994introduction} descrito en la Ecuación (13), ullamcorper suscipit lobortis nisl ut aliquip ex
ea commodo consequat rápidamente porque son puntos que poseen al menos una variable adimensional compleja, es decir, un campo $\Psi$ con componente imaginaria.}


\section*{Resultados}

Deben presentarse en secuencia lógica, con las tablas y las figuras acompañadas de una explicación y análisis de estas.

\section*{Discusión}

Discusión y limitaciones de los hallazgos y comparar las observaciones relevantes con las de otros estudios.
Los resultados y la discusión pueden presentarse también en forma combinada, en cuyo caso la sección se denominará Resultados y Discusión.


\begin{figure}[H]
    \centering
    \includegraphics[scale=0.4]{Imagenes/85-anos.jpg}
    \caption{(Las fotografías, gráficos, dibujos y esquemas se denominan todos como Figuras, las cuales se enumeran según el orden de aparición. Si se trata de micro-fotografías debe indicarse el aumento utilizado. Las personas, o sus nombres, no
deben ser identificables en las fotografías, pero se requiere permiso escrito para su publicación)}
    \label{fig:my_label}
\end{figure}


\begin{table}[H]
\caption{Las tablas y cuadros se denominarán tablas y llevarán numeración arábiga de acuerdo con el orden de aparición. El título correspondiente se coloca en la parte superior de la tabla y las notas en la inferior. Los símbolos para unidades deben aparecer en el encabezamiento de las columnas. Si es necesario exponer alguna información, utilizar las letras minúsculas del alfabeto según orden de aparición en la respectiva tabla: a, b, c, d, e,
etc.}
\begin{tabular}{llll}
\hline
\textbf{Letra} & \textbf{Tamaño} & \textbf{Letra} & \textbf{Justificación} \\ \hline
\rowcolor[HTML]{EFEFEF} 
Titulo         & 16pt            & Arial          & Centrada               \\ \hline
Autores        & 10pt            & Arial          & Centrada               \\ \hline
\rowcolor[HTML]{EFEFEF} 
Texto          & 10pt            & Arial          & Justificada            \\ \hline
Ecuación       & 10pt            & Arial          & Justificada            \\ \hline
\rowcolor[HTML]{EFEFEF} 
Figura         & 10pt            & Arial          & Justificada            \\ \hline
Tabla          & 10pt            & Arial          & Justificada            \\ \hline
\end{tabular}
\label{Tab:Ejemplo_Tabla}
\end{table}


\section*{Conclusions}

Deberán responder al objetivo de la investigación reportada e incluir el planteamiento de objetivos futuros a seguir en la línea de investigación.

\section*{Agradecimientos}

Podrán ser expresados a personas o entidades que hayan ayudado a en la realización del trabajo. Se mencionarán las personas y el tipo de ayuda aportada, así como las entidades, instituciones o fundaciones que hayan financiado o suministrado materiales
En caso de proyectos de investigación financiados será suficiente con mencionar el código de registro y la entidad, institución o fundación que lo apoyó económicamente.

\section*{Contribución de los autores}

Describe con precisión las contribuciones individuales de cada autor al trabajo, señalando el nombre del autor con siglas (Ejemplo: Elizabeth Castañeda debe ir EC: Contribución...).
Las contribuciones de todos los autores deben estar definidas en el manuscrito enviado y se publicarán en el artículo final.


\section*{Conflicto de intereses}

Si los autores o coautores, editores o revisores (evaluadores) tienen algún conflicto de intereses que pueda de alguna manera influir en la transparencia u objetividad en el proceso de revisión por pares y publicación, tal conflicto debe ser declarado.
Los conflictos de interés más evidentes son los llamados conflictos financieros, resultado de relaciones entre empleador-empleado, financiaciones y patentes, así como también honorarios,
consultorías, investigación pagada por compañías que dependen de los resultados investigativos, entre otros.


%%%% Bibliografía
\bibliography{Bibliografia.bib}
\bibliographystyle{apacite} %%No cambiar el estilo de la bibliografía, el formato es propio de la revista

\end{document}
